\documentclass{article}
\usepackage[utf8]{inputenc}
\usepackage{amsmath}
\usepackage{amssymb}
\usepackage{hyperref}
\usepackage{listings}
\usepackage{geometry}
\geometry{a4paper, margin=1in}

\title{Technical Report: Real-time Bethe-Salpeter Equation (BSE) Implementation in paraHFNAMD}
\author{Gemini CLI}
\date{March 2, 2026}

\begin{document}

\maketitle

\section{Introduction}
The \texttt{paraHFNAMD} package implements a real-time Bethe-Salpeter Equation (BSE) solver designed to calculate exciton properties based on ground-state GW calculations. This report details the implementation of the BSE module, the functions involved, and their assembly into the overall workflow.

\section{Core Classes and Data Structures}
The BSE implementation is primarily housed in the \texttt{src/bse.cpp} and \texttt{src/bse.h} files, centered around the \texttt{excitonclass}. It interacts with several other key components:
\begin{itemize}
    \item \textbf{\texttt{waveclass}}: Manages electronic wavefunctions, quasiparticle energies, and PAW (Projector Augmented Wave) information.
    \item \textbf{\texttt{wpotclass}}: Handles the screened Coulomb potential ($W$) obtained from a prior GW calculation.
    \item \textbf{\texttt{socclass}}: Handles spin-orbit coupling effects if requested.
\end{itemize}

\section{BSE Hamiltonian Construction}
The BSE Hamiltonian $H^{BSE}$ is constructed in the basis of electron-hole pairs ($ck, vk'$). In the Tamm-Dancoff approximation, the matrix elements are:
\begin{equation}
H_{(ck,vk'), (c'k'',v'k''')} = (\epsilon_{ck} - \epsilon_{vk})\delta_{cc'}\delta_{vv'}\delta_{kk''}\delta_{k'k'''} + K_{(ck,vk'), (c'k'',v'k''')}
\end{equation}
where $K = K^d + \kappa K^x$ is the interaction kernel.

\section{Practical Implementation of $W_{GG'}(q)$}
The screened Coulomb potential $W_{GG'}(q)$ is a central component of the BSE kernel. In \texttt{paraHFNAMD}, this is handled by the \texttt{wpotclass}.

\subsection{Reading $W$ from Files}
The potential is typically read from binary files (e.g., \texttt{Wxxxx.tmp} or \texttt{WFULLxxxx.tmp}) generated by VASP's GW implementation.
\begin{itemize}
    \item \textbf{Head and Wing}: The $q \to 0$ limits ($W_{00}$ and $W_{0G}$) are read from \texttt{W0001.tmp} in \texttt{ReadHeadWing()}. The head represents the average of the diagonal dielectric tensor elements.
    \item \textbf{Full Matrix}: The \texttt{WFULLxxxx.tmp} files contain the full $W_{GG'}(q)$ matrices for each k-point in the irreducible Brillouin zone.
    \item \textbf{Memory-Efficient Reading}: To handle large matrices, \texttt{ReadOneWFullSub()} reads the matrix in blocks corresponding to irreducible k-points, minimizing the instantaneous memory footprint.
\end{itemize}

\subsection{Reciprocal Space Mapping ($Q+G \to q+g$)}
GW calculations often use a primitive cell, while BSE may use a supercell. The mapping is performed in \texttt{QpG2qpg()}:
\begin{enumerate}
    \item A transformation matrix $M$ is computed to relate supercell and primitive cell lattice vectors.
    \item For each supercell $Q+G$ vector, the corresponding primitive cell $q+g$ vector is identified.
    \item Symmetry operations (rotations, time-reversal) are applied using \texttt{Doublex012toqg()} to map any $q$-point back to the irreducible Brillouin zone where $W$ is stored.
\end{enumerate}

\section{Practical Calculation of Interaction Kernels}

\subsection{Direct Interaction ($K^d$)}
The calculation of $K^d$ is implemented in \texttt{DirectTerms()}.
\begin{itemize}
    \item \textbf{Data Distribution}: The calculation is parallelized over k-point pairs $(k, k')$ using MPI. Each process handles a subset of pairs.
    \item \textbf{Matrix Operations}: For each pair, the code:
    \begin{enumerate}
        \item Computes density matrices $M_{cc'}(k, q, G)$ and $M_{vv'}(k, q, G')$ using PAW-augmented wavefunctions.
        \item Performs the double contraction: $K^d = \alpha \times \text{trans}(M_{cc'}) \times W \times \text{conj}(M_{vv'})$.
        \item Uses high-performance BLAS \texttt{Zgemm} for the matrix multiplications.
    \end{enumerate}
    \item \textbf{Singularity Handling}: For $q=0$, the head and wing contributions are added separately to handle the Coulomb singularity.
\end{itemize}

\subsection{Exchange Interaction ($K^x$)}
The exchange term is implemented in \texttt{ExchangeTerms()}.
\begin{itemize}
    \item \textbf{Density Matrix}: It first computes the $q=0$ density matrix $M_{cv}(k, 0, G)$ using \texttt{DensityMatrixKCV()}.
    \item \textbf{ScaLAPACK Integration}: Because the resulting BSE matrix is globally distributed in a 2D block-cyclic manner, the exchange term (which involves all excitonic transitions) is added using ScaLAPACK's \texttt{Pzgemm}.
    \item \textbf{Scaling}: Each term is scaled by the bare Coulomb potential $v(G) = 4\pi/|G|^2$ and the crystal volume $\Omega$.
\end{itemize}

\section{Numerical Diagonalization and Phase Alignment}
After the Hamiltonian is assembled, \texttt{XEnergies()} performs the following:
\begin{enumerate}
    \item \textbf{Diagonalization}: Uses \texttt{Pzheev} (ScaLAPACK) to find eigenvalues and eigenvectors.
    \item \textbf{Phase Continuity}: In a real-time trajectory, the phase of eigenvectors can jump between steps. The code aligns the current eigenvectors $A(t)$ with the previous ones $A(t-dt)$ by calculating the overlap $\langle A(t-dt) | A(t) \rangle$ and applying a phase correction. This is essential for stable TDM and non-adiabatic coupling calculations.
    \item \textbf{Exciton TDM}: Transition dipole moments are computed in \texttt{ExcitonTDM()} by weighting the momentum matrix elements $\mathbf{P}_{cvk}$ with the exciton coefficients $A^S_{cvk}$.
\end{enumerate}

\section{Summary of Functional Assembly}
The \texttt{OnlyBSECalc()} driver coordinates the following sequence for each time step:
\begin{enumerate}
    \item \textbf{Wavefunction Load}: \texttt{WVCBasicProcess()} reads electronic states.
    \item \textbf{Mapping}: \texttt{wpotclass} maps GW potential to the current basis.
    \item \textbf{Kernel Construction}: \texttt{DirectTerms()} and \texttt{ExchangeTerms()} build the interaction matrices.
    \item \textbf{Matrix Assembly}: \texttt{ScTerms2BSEMat()} and \texttt{ExTerms2BSEMat()} scatter local contributions into the global block-cyclic BSE Hamiltonian.
    \item \textbf{Solution}: \texttt{XEnergies()} diagonalizes the Hamiltonian and extracts exciton properties.
\end{enumerate}

\end{document}
